\section{Finite resolutions for the analytic groups}
\label{sec:finite-resolutions}

The completed coefficient calculation uses continuous cohomology with
noncompact valuation coefficients. We first give the finite resolution
that computes those cochains and bounds their cohomological degrees.
This tool depends only on the analytic groups and coefficient topology.
A linearly topologized \(\Fp\)-space is understood to be separated; an
action has \emph{invariant neighborhoods} when invariant open linear
subspaces form a neighborhood basis of zero.

\begin{lemma}[Finite analytic resolutions]
\label{lem:finite-resolution}
Let $G$ be a compact $p$-adic analytic group with no element of order $p$,
and put $d=\dim G$.  The trivial module $\Fp$ admits a length-$d$
resolution $Q_\bullet$ by finitely generated projective
$\Fp[[G]]$-modules, with their compact topologies.  If $M$ is a complete
linearly topologized $\Fp$-representation with invariant neighborhoods,
then
\[
 \Hom_{\Fp[[G]],\mathrm{cts}}(Q_\bullet,M)
 \simeq \Ccts(G,M)
\]
by continuous chain-homotopy equivalences.  Each term on the left is a
continuous direct summand of a finite power of $M$.

For $G=H$ the dimension is eight, and the reversed group-ring dual of
$Q_\bullet$ resolves the trivial right $\Fp[[H]]$-module.  In particular,
continuous $H$-cohomology of these coefficients vanishes above degree
eight, and finite coefficient modules have finite cohomology.
\end{lemma}

\begin{proof}
We specify the two structural inputs.  A completed group algebra over
$\Zp$, or over its residue field, is Noetherian; it has finite global
dimension when the compact analytic group has no element of order $p$
\cite[Section~3.3, Proposition~(d),(e)]{AW06}.  Also
\begin{equation}
 \label{eq:analytic-dualizing-cohomology}
 H^a_{\mathrm{cts}}(G,\Zp[[G]])=
 \begin{cases}
  \mathcal D_G,&a=d,\\
  0,&a\ne d,
 \end{cases}
\end{equation}
where $\mathcal D_G$ is free of rank one over $\Zp$, with the right
action given by the top exterior power of the dual adjoint
representation \cite[Proposition~4.16, Remark~4.23, and
Proposition~4.40]{BGHS22}.

Set $\widetilde\Lambda=\Zp[[G]]$ and $\Lambda=\Fp[[G]]$.
Noetherianity and finite global dimension give a finite resolution of
$\Zp$ by finitely generated projective $\widetilde\Lambda$-modules.
Give every term its topology as a summand of a finite power of
$\widetilde\Lambda$.  The differentials are continuous and their images
are compact, hence closed.  All syzygies are $p$-torsion-free, so
reduction modulo $p$ remains exact and gives a finite projective
$\Lambda$-resolution of $\Fp$.

Here the integral projective resolution computes
\eqref{eq:analytic-dualizing-cohomology} as continuous cohomology.
Indeed, a completed free $\Zp$-module on a profinite set is compact and
$p$-torsion-free.  Its Pontryagin dual is a divisible discrete
$p$-primary group, and hence injective.  It is therefore projective in
compact $\Zp$-modules.  Induction to $\widetilde\Lambda$ shows that the
completed bar terms are projective in compact
$\widetilde\Lambda$-modules.  The usual projective comparison maps and
homotopies can consequently be chosen continuous.  Dualizing the finite
integral resolution and then reducing modulo $p$ gives cohomology
$\mathcal D_G/p$ in degree $d$ alone: both the integral terms and the
unique integral cohomology module are $p$-torsion-free.  Finite
projectivity identifies this reduction with the group-ring dual of the
reduced resolution.  If its last degree is $n>d$, the last differential
of that dual surjects onto a projective module.  Split this surjection
and dualize back to remove a contractible pair.  Repetition gives a
resolution of length $d$.

We explain why the same compact resolution works for the possibly
noncompact target $M$.  For a profinite space $Z$, there is a natural
identification
\begin{equation}
 \label{eq:completed-free-continuous-functions}
 C_{\mathrm{cts}}(Z,M)
 =\Hom_{\Fp,\mathrm{cts}}(\Fp[[Z]],M).
\end{equation}
Modulo each open linear subspace of $M$, a continuous function has
finite image and factors through a finite clopen partition of $Z$.
Its linear extension therefore exists in that quotient.  The extensions
are compatible, and completeness gives the required map into $M$.
This also proves uniqueness.  With invariant neighborhoods the
$G$-action extends continuously to $\Lambda$.

The completed $\Fp$-bar terms are projective in compact
$\Lambda$-modules.  To lift a map from their profinite generators across
a compact-module surjection, first choose a continuous linear section
of the underlying compact $\Fp$-spaces.  Such a section exists by
dualizing and splitting an injection of discrete vector spaces.
Extend the lifted generator map $\Lambda$-linearly using
\eqref{eq:completed-free-continuous-functions}.  Thus the completed bar
resolution and $Q_\bullet$ are continuously chain-homotopy equivalent.
Apply continuous $\Lambda$-linear Hom into $M$.  Continuity of the
resulting maps and homotopies for the uniform cochain topologies is
checked modulo invariant open subspaces of $M$.  Finally, finite
projectivity makes each Hom term a continuous summand of $M^r$ for
some finite $r$.

For $H$, an element of order three would embed
$\mathbb Q_3(\zeta_3)$ in the division algebra $D_3$.  The degree of a
commutative subfield divides the degree three of this central division
algebra, whereas $[\mathbb Q_3(\zeta_3):\mathbb Q_3]=2$.  Thus $H$ has
no element of order three.  Its Lie algebra is the kernel of reduced
trace on $D_3$, so it has dimension eight.  Over a splitting field,
conjugation on $D_3$ becomes conjugation on $M_3$ and has determinant
one.  The decomposition
\[
 D_3=\Qp\cdot1\oplus\ker(\operatorname{Trd})
\]
is valid over $\Qp$, including at $p=3$, and the first summand has
trivial action.  The adjoint determinant on $\Lie(H)$ is therefore one.
Consequently $\mathcal D_H$ is the trivial line, which proves the
assertion about the reversed dual.  The remaining statements follow
by applying the finite resolution.
\end{proof}

For the torsion-free uniform open subgroups of \(J\), the same lemma
applies with dimension four. For \(H\), it supplies the dimension-eight
resolution and duality orientation needed later in the finite-stage
argument. We now turn to the completed geometric coefficients.
